\section*{Hoofdstuk \ref{ch:b3:complex-mechanisms} --- Reactiemechanismen van complexen}

\begin{solution}{exo:b3:complex-mechanisms:1}
Inert: $\ce{[Cr(H2O)6]^{3+}}$ ($d^3$) en $\ce{[Co(NH3)6]^{3+}}$ (laagspin-$d^6$), grote
ligandveldactiveringsenergieën. \omterm{def:b3:complex-mechanisms:labile}{Labiel}: $\ce{[Cr(H2O)6]^{2+}}$ (hoogspin-$d^4$) en $\ce{[Cu(H2O)6]^{2+}}$ ($d^9$),
door Jahn-Teller vervormd met lange, zwakke axiale bindingen; $\ce{[Zn(H2O)6]^{2+}}$ ($d^{10}$), geen
ligandveldbarrière.
\end{solution}

\begin{solution}{exo:b3:complex-mechanisms:2}
Als $k_2[\mathrm Y] \gg k_{-1}[\mathrm X]$ is de noemer $k_2[\mathrm Y]$ en $v = k_1[\mathrm{ML_5X}]$; als $k_{-1}[\mathrm X] \gg k_2[\mathrm Y]$ is hij
$k_{-1}[\mathrm X]$ en volgt de tweede vorm (een voorevenwicht dat door het uittredende
ligand geremd wordt).
\end{solution}

\begin{solution}{exo:b3:complex-mechanisms:3}
D: $\Delta V^{\ddagger} > 0$ en $\Delta^{\ddagger}S^\circ > 0$ (een ligand vertrekt). A: beide negatief (een ligand wordt in
de overgangstoestand gebonden).
\end{solution}

\begin{solution}{exo:b3:complex-mechanisms:4}
$cis$-$\ce{[PtCl2(NH3)2]}$ uit $\ce{[PtCl4]^{2-}}$; $trans$-$\ce{[PtCl2(NH3)2]}$ uit $\ce{[Pt(NH3)4]^{2+}}$ (het
\omterm{def:b3:complex-mechanisms:trans-effect}{trans-effect} van \ce{Cl-} is groter dan dat van \ce{NH3}).
\end{solution}

\begin{solution}{exo:b3:complex-mechanisms:5}
$k_{\mathrm{obs}} = 2.0 \times \num{3.0e4} \times 0.10/1.2 = \qty{5.0e3}{s^{-1}}$; bij \qty{1.0}{mol/L} is het $\num{6.0e4}/3.0 = \qty{2.0e4}{s^{-1}}$; limiet
$k_i = \qty{3.0e4}{s^{-1}}$.
\end{solution}

\begin{solution}{exo:b3:complex-mechanisms:6}
$\Delta V^{\ddagger} = -RT\ln2/\Delta p = -8.314 \times 298 \times 0.693/\num{99.9e6} = \qty{-1.7e-5}{m^3.mol^{-1}} = \qty{-17}{cm^3.mol^{-1}}$:
associatief.
\end{solution}

\begin{solution}{exo:b3:complex-mechanisms:7}
Het fosfine, een sterke $\sigma$-donor en $\pi$-acceptor, verzwakt de binding die er trans tegenover staat (\omterm{def:b3:complex-mechanisms:trans-effect}{trans-invloed},
langere binding) en maakt haar \omterm{def:b3:complex-mechanisms:labile}{labiel} (\omterm{def:b3:complex-mechanisms:trans-effect}{trans-effect}): dat chloride wordt
als eerste vervangen.
\end{solution}

\begin{solution}{exo:b3:complex-mechanisms:8}
Het brugligand komt terecht op het geoxideerde metaal; de snelheid hangt
sterk af van de aard van de mogelijke brug; er wordt een tweekernig
intermediair gedetecteerd; en de snelheid kan niet groter zijn dan die
van de substitutie die nodig is om de brug op de labiele partner te
vormen.
\end{solution}

\begin{solution}{exo:b3:complex-mechanisms:9}
$(\lambda + \Delta G^\circ)^2/4\lambda$ met $4\lambda = 320$: 20.0, 11.3, 0 en \qty{5.0}{kJ/mol} (de laatste in het
geïnverteerde gebied).
\end{solution}

\begin{solution}{exo:b3:complex-mechanisms:10}
$k_{12} = \sqrt{4.0 \times \num{2.0e5} \times \num{1.0e4}} = \qty{8.9e4}{L.mol^{-1}.s^{-1}}$.
\end{solution}

\begin{solution}{exo:b3:complex-mechanisms:11}
Hoogspin-$d^5$: 0; $d^7$: $1.33\,Dq$; $d^8$: $2.0\,Dq$. Verwachte snelheden van wateruitwisseling: \ce{Mn^{2+}} $>$ \ce{Co^{2+}} $>$
\ce{Ni^{2+}}, de waargenomen volgorde.
\end{solution}

\begin{solution}{exo:b3:complex-mechanisms:12}
Bij zeer exergonische reacties is de intrinsieke snelheid groter dan de
snelheid waarmee de ionen elkaar ontmoeten: de waargenomen constante
blijft op de diffusielimiet en de daling blijft verborgen. Met donor en
acceptor op een vaste afstand in één molecuul is er geen diffusiestap,
en een reeks acceptoren met toenemende drijvende kracht toonde een
snelheid die steeg, door een maximum ging en daalde.
\end{solution}

\begin{solution}{pb:b3:complex-mechanisms:1}
\textbf{1.} $5d^8$.
\textbf{2.} Het sterke veld van een $5d$-ion maakt de $d_{x^2-y^2}$-orbitaal, die naar vier
liganden wijst, zeer hoog: acht elektronen paren in de vier lagere
orbitalen.
\textbf{3.} Het complex heeft 16 valentie-elektronen en een vrije axiale
plaats: het binnenkomende ligand bindt eerst, via een vijfgecoördineerde
overgangstoestand.
\textbf{4.} $k_{\mathrm{obs}} = k_1 + k_2[\mathrm Y]$: een oplosmiddelweg (aanval van het oplosmiddel, daarna snelle vervanging door Y)
en een rechtstreekse weg.
\textbf{5.} Substitutie duurt uren (deel III): inert op de schaal van
minuten, niet op die van een dag.
\textbf{6.} Cisplatine, omdat de tweede \ce{NH3} een chloride vervangt dat trans staat ten opzichte van chloride
(\omterm{def:b3:complex-mechanisms:trans-effect}{trans-effect} $\ce{Cl-} > \ce{NH3}$); er ontstaan ook nevenproducten.
\textbf{7.} Jodide heeft een groter \omterm{def:b3:complex-mechanisms:trans-effect}{trans-effect} dan chloride, wat het
sturende effect zuiverder en sneller maakt.
\textbf{8.} In $\ce{[PtI3(NH3)]-}$ wordt het jodide trans ten opzichte van jodide \omterm{def:b3:complex-mechanisms:labile}{labiel} gemaakt, niet dat trans ten opzichte van
\ce{NH3}: de tweede ammoniak komt cis ten opzichte van de eerste binnen.
\textbf{9.} Zilverjodide slaat neer en $cis$-$\ce{[Pt(H2O)2(NH3)2]^{2+}}$ blijft in oplossing.
\textbf{10.} Chloride vervangt de twee watermoleculen op hun eigen
plaatsen; de ammines worden niet aangeraakt.
\textbf{11.} Uit $\ce{[Pt(NH3)4]^{2+}}$ en twee chloriden.
\textbf{12.} Zijn twee reactieve posities liggen tegenover elkaar: het kan
geen twee naburige basen van dezelfde DNA-streng binden, de laesie
waarmee cisplatine werkt.
\textbf{13.} $\ln2/\num{9.0e-5} = \qty{7.7e3}{s} = \qty{2.1}{h}$.
\textbf{14.} $1 - \eu^{-\num{9.0e-5} \times 7200} = 0.48$.
\textbf{15.} Water is een veel betere uittredende groep dan chloride: het
aquacomplex wordt snel vervangen door het stikstofatoom van een DNA-base.
\textbf{16.} Associatief, zoals alle substituties van platina(II): een
negatief \omterm{def:b3:complex-mechanisms:activation-volume}{activeringsvolume}.
\textbf{17.} De hoge chlorideconcentratie duwt het evenwicht terug naar de
dichlorovorm, zodat het geneesmiddel intact blijft tot het de cellen
bereikt.
\textbf{18.} $f = K/(K + [\ce{Cl-}])$.
\textbf{19.} $\num{3.0e-3}/0.103 = 0.029$.
\textbf{20.} $\num{3.0e-3}/0.0070 = 0.43$.
\textbf{21.} 15.
\textbf{22.} Alleen binnen cellen bestaat een aanzienlijke fractie als het
reactieve aquacomplex.
\textbf{23.} Zwavelliganden, zoals glutathion en het cysteïne en methionine
van eiwitten, die het zachte platina sterk binden.
\textbf{24.} Ammoniak bindt platina sterk en staat trans ten opzichte van
chloride, een ligand met een klein \omterm{def:b3:complex-mechanisms:trans-effect}{trans-effect}: de amminebindingen
worden niet \omterm{def:b3:complex-mechanisms:labile}{labiel} gemaakt.
\textbf{25.} \textbf{Binnen een cel is ongeveer 15 keer meer van het
geneesmiddel in de reactieve aquavorm aanwezig dan in plasma.}
\end{solution}
